\section*{Chapter \ref{ch:b2:plant-life-cycles} --- Life Cycles and Reproduction of Land Plants}

\begin{solution}{exo:b2:plant-life-cycles:1}
\omterm{def:b2:plant-life-cycles:cycles}{Gametophyte}: the \omterm{def:b2:meiosis-heredity:meiosis}{haploid} \omterm{def:b2:unicellular-diversity:unicellular}{multicellular} generation, grown by mitosis
from a spore, that makes gametes by mitosis. \omterm{def:b2:plant-life-cycles:cycles}{Sporophyte}: the \omterm{def:b2:meiosis-heredity:meiosis}{diploid}
generation, grown by mitosis from a zygote, that makes spores by
\omterm{def:b2:meiosis-heredity:meiosis}{meiosis}. Spore: a \omterm{def:b2:meiosis-heredity:meiosis}{haploid} cell made by \omterm{def:b2:meiosis-heredity:meiosis}{meiosis} that grows without
fusing. Gamete: a \omterm{def:b2:meiosis-heredity:meiosis}{haploid} cell made by mitosis (in plants) that must
fuse with another.
\end{solution}

\begin{solution}{exo:b2:plant-life-cycles:2}
Leaf cells are $2n = 1260$: spore 630, \omterm{def:b2:plant-life-cycles:fern}{prothallus} cell 630, sperm
630, egg 630, zygote 1260.
\end{solution}

\begin{solution}{exo:b2:plant-life-cycles:3}
Green cushion: \omterm{def:b2:plant-life-cycles:cycles}{gametophyte}. Stalk and capsule: \omterm{def:b2:plant-life-cycles:cycles}{sporophyte}. \omterm{def:b2:plant-life-cycles:fern}{Fern}
frond: \omterm{def:b2:plant-life-cycles:cycles}{sporophyte}. \omterm{def:b2:plant-life-cycles:fern}{Prothallus}: \omterm{def:b2:plant-life-cycles:cycles}{gametophyte}. \omterm{def:b2:plant-life-cycles:heterospory}{Pollen} grain: male
\omterm{def:b2:plant-life-cycles:cycles}{gametophyte}. Seed's food store (in a pine): female \omterm{def:b2:plant-life-cycles:cycles}{gametophyte}.
Seed's embryo: the next \omterm{def:b2:plant-life-cycles:cycles}{sporophyte}.
\end{solution}

\begin{solution}{exo:b2:plant-life-cycles:4}
Haplontic (\emph{Chlamydomonas}): \omterm{def:b2:meiosis-heredity:meiosis}{meiosis} in the zygote, at once.
Diplontic (animals, \emph{Fucus}): \omterm{def:b2:meiosis-heredity:meiosis}{meiosis} makes the gametes.
Haplo-diplontic (mosses, \omterm{def:b2:plant-life-cycles:fern}{ferns}, seed plants, many algae): \omterm{def:b2:meiosis-heredity:meiosis}{meiosis}
makes spores in the \omterm{def:b2:plant-life-cycles:cycles}{sporophyte}'s sporangia.
\end{solution}

\begin{solution}{exo:b2:plant-life-cycles:5}
$0.03/10^{-4} = \qty{300}{s}$, five minutes. In \qty{20}{min},
\qty{12}{cm}. Fertilisation works only between plants a few
centimetres apart, so mosses grow in dense cushions with both sexes
(or both organs) within reach.
\end{solution}

\begin{solution}{exo:b2:plant-life-cycles:6}
$v_s = 2\times(1.5\times 10^{-5})^2\times 1099\times 9.81/(9\times
1.8\times 10^{-5}) = \qty{0.030}{m/s}$, \qty{3}{cm/s}; distance $uh/v_s
= 2\times 0.5/0.03 = \qty{33}{m}$. Seed: $2\times 20/0.8 =
\qty{50}{m}$ --- a heavier unit shed from higher with a wing goes
about as far.
\end{solution}

\begin{solution}{exo:b2:plant-life-cycles:7}
$200\times 40\times 64 = \num{512000}$ spores per frond; $5.1$
prothalli; $0.1$ young \omterm{def:b2:plant-life-cycles:cycles}{sporophyte} per frond.
\end{solution}

\begin{solution}{exo:b2:plant-life-cycles:8}
A seed is a retained \omterm{def:b2:plant-life-cycles:heterospory}{ovule}: the female \omterm{def:b2:plant-life-cycles:cycles}{gametophyte} must be kept on
the parent, enclosed, and fed, while the male \omterm{def:b2:plant-life-cycles:cycles}{gametophyte} must
travel to it. That needs two kinds of spore with two fates. A
homosporous plant's one spore must be shed to grow into a free
bisexual \omterm{def:b2:plant-life-cycles:cycles}{gametophyte}; retaining it would retain the male function as
well and force self-fertilisation on every plant, and there would be
nothing to travel.
\end{solution}

\begin{solution}{exo:b2:plant-life-cycles:9}
Spore: $\tfrac43\pi(1.5\times 10^{-5})^3\times 1100 =
\qty{1.6e-11}{kg}$, \qty{16}{ng}; energy $0.5\times 1.6\times
10^{-8}\times 2\times 10^{4} = \qty{1.6e-4}{J}$. Seed: \qty{5}{mg},
$0.5\times 5\times 10^{-3}\times 2\times 10^{4} = \qty{50}{J}$ ---
$3\times 10^{5}$ times more. The spore can grow one small
photosynthetic scale and nothing before light; the seed can grow a
root and a shoot in the dark, wait through a season, and survive
being eaten or dried.
\end{solution}

\begin{solution}{exo:b2:plant-life-cycles:10}
$10^{11}/(10^{6}\times 20) = \num{5000}$ grains per cubic metre; flux
$5000\times 10^{-7}\times 2 = \qty{1e-3}{}$ grains per second; about
\qty{1000}{s}, a quarter of an hour, per grain. The cone must be open
and receptive for hours to days, and precisely when the male cones
shed --- pollination is timed to the week.
\end{solution}

\begin{solution}{exo:b2:plant-life-cycles:11}
Diploidy masks \omterm{def:b2:meiosis-heredity:vocabulary}{recessive} deleterious \omterm{def:b2:genome-diversification:mutation}{mutations}; a \omterm{def:b2:meiosis-heredity:meiosis}{diploid} body with
vascular tissue can be large and long-lived, and height helps both
light capture and \omterm{thm:b2:plant-life-cycles:dispersal}{spore dispersal}; spores made high on a \omterm{def:b2:plant-life-cycles:cycles}{sporophyte}
disperse in air, whereas gametes must swim. The \omterm{def:b2:meiosis-heredity:meiosis}{haploid} generation
persists because \omterm{def:b2:meiosis-heredity:meiosis}{meiosis} and fertilisation require a \omterm{def:b2:meiosis-heredity:meiosis}{haploid} phase,
and the \omterm{def:b2:plant-life-cycles:heterospory}{pollen} grain and embryo sac are the minimal bodies that
carry gametes to each other and feed the embryo.
\end{solution}

\begin{solution}{exo:b2:plant-life-cycles:12}
By culturing spores and following them under the microscope he could
establish that a spore grows into a small body bearing sex organs,
that the embryo arises from the fertilised egg inside the
\omterm{def:b2:plant-life-cycles:moss}{archegonium}, that the spore-bearing plant grows from it, and that
the same organs exist in reduced form in the conifer \omterm{def:b2:plant-life-cycles:heterospory}{ovule}: the
alternation as a sequence of bodies and organs. He could not know
what distinguished the two generations at the level of the cell.
Chromosome counts showed that the two bodies differ by a factor two
in chromosome number, that \omterm{def:b2:meiosis-heredity:meiosis}{meiosis} sits at spore formation and
fertilisation at the egg, and so explained why the alternation is
obligatory.
\end{solution}

\begin{solution}{pb:b2:plant-life-cycles:1}
\textbf{1.} $10\times 200\times 40\times 64 = 5.1\times 10^{6}$
spores.
\textbf{2.} $\tfrac43\pi(1.5\times 10^{-5})^3\times 1100 =
\qty{1.6e-11}{kg}$ per spore; crop $8\times 10^{-5}$ kg, \qty{80}{mg}.
\textbf{3.} $v_s = 2\times 2.25\times 10^{-10}\times 1099\times
9.81/(1.62\times 10^{-4}) = \qty{0.030}{m/s}$.
\textbf{4.} $x = 1.5\times 0.5/0.030 = \qty{25}{m}$; disc of
$\pi\times 25^2 = \qty{2000}{m^2}$; about \num{2600} spores per
square metre.
\textbf{5.} Updrafts of a few centimetres per second exceed $v_s$, so
a spore caught in turbulence stays aloft for hours and travels
hundreds of kilometres; a single spore founds a population if its
\omterm{def:b2:plant-life-cycles:fern}{prothallus} can self-fertilise (it bears both organs), which is why
remote islands have rich \omterm{def:b2:plant-life-cycles:fern}{fern} floras.
\textbf{6.} Spore: $0.5\times 1.6\times 10^{-8}\times 2\times 10^{4} =
\qty{1.6e-4}{J}$; crop: \qty{820}{J}; one pine seed: \qty{50}{J} ---
the whole spore crop equals sixteen seeds.
\textbf{7.} $5.1\times 10^{6}/2\times 10^{4} = 256$ prothalli.
\textbf{8.} $1 - 0.8^{6} = 1 - 0.26 = 0.74$.
\textbf{9.} $256\times 0.74\times 0.5\times 0.02 = 1.9$ adult \omterm{def:b2:plant-life-cycles:fern}{ferns}
per season.
\textbf{10.} $30\times 1.9 \approx 57$: far more than the one
replacement of a stable population, so the survival figures are
optimistic --- in a full habitat most young \omterm{def:b2:plant-life-cycles:cycles}{sporophytes} die of
crowding and shade.
\textbf{11.} $10^{-3}/10^{-4} = \qty{10}{s}$. Many \omterm{def:b2:plant-life-cycles:fern}{ferns} mature
antheridia and archegonia at different times, or release a hormone
(antheridiogen) that makes neighbouring young prothalli male, so
that sperm come from another individual.
\textbf{12.} It is one cell thick, rootless, unprotected, needs
liquid water for fertilisation and lives a few weeks: nearly all the
mortality of the cycle falls on it (one spore in $2\times 10^{4}$
becomes one; a quarter of those never see a wet night).
\textbf{13.} $v_s = 2\times 6.25\times 10^{-10}\times 599\times
9.81/(1.62\times 10^{-4}) = \qty{0.045}{m/s}$.
\textbf{14.} $3\times 15/0.045 = \qty{1000}{m}$.
\textbf{15.} $10^{11}/(2\times 10^{7}) = \num{5000}$ per cubic metre.
\textbf{16.} $5000\times 10^{-7}\times 2 = \qty{1e-3}{s^{-1}}$: 3.6 per
hour; the first arrives after about \qty{1000}{s}, a quarter of an
hour.
\textbf{17.} Five kilometres downwind the cloud has spread sideways
and upward and most grains have settled (a \qty{1}{km} range per
\qty{15}{m} of height), so the concentration is orders of magnitude
lower and an \omterm{def:b2:plant-life-cycles:heterospory}{ovule} may wait days for a grain; the tree's own \omterm{def:b2:plant-life-cycles:heterospory}{pollen}
mostly gives selfed seeds, which abort. Wind pollination needs a
dense cloud, hence stands.
\textbf{18.} $10^{14}$ grains for $2\times 10^{6}$ \omterm{def:b2:plant-life-cycles:heterospory}{ovules}: $5\times
10^{7}$ grains per \omterm{def:b2:plant-life-cycles:heterospory}{ovule}.
\textbf{19.} $2\times 10^{6}$ \omterm{def:b2:plant-life-cycles:heterospory}{ovules} $\times 0.6\times 0.8 = 9.6\times
10^{5}$ seeds; \qty{4.8}{kg}; \qty{48}{MJ}.
\textbf{20.} $3\times 20/0.8 = \qty{75}{m}$: the \omterm{def:b2:plant-life-cycles:heterospory}{pollen} goes a
kilometre, the seed a few tree-heights; the genes travel by \omterm{def:b2:plant-life-cycles:heterospory}{pollen},
the plant by seed.
\textbf{21.} $9.6\times 10^{5}\times 0.05\times 0.01 = 480$ adult
offspring, against the \omterm{def:b2:plant-life-cycles:fern}{fern}'s 57.
\textbf{22.} Pine: $4.8\times 10^{7}/480 = \qty{100}{kJ}$ per adult
offspring; \omterm{def:b2:plant-life-cycles:fern}{fern}: $820/1.9 = \qty{430}{J}$ per adult offspring, two
hundred times less.
\textbf{23.} Both are half dry matter at \qty{20}{kJ/g}, so at a
basal rate of, say, \qty{0.5}{mW} per gram dry both last $10^{4}\,
\text{J/g}/(5\times 10^{-4}\,\text{W/g}) = 2\times 10^{7}$ s, about
eight months: the duration of dormancy is not what the seed's energy
buys. It buys what happens after germination --- a root and a shoot
built in the dark from reserves, weeks before the seedling feeds
itself --- and, with the coat, protection meanwhile; the spore's
\qty{1.6e-4}{J} builds a few cells that must photosynthesise at once.
\textbf{24.} Seed: fertilisation without water; a provisioned,
protected, dormant embryo that establishes in dry or seasonal
places; dispersal by wings, fruits and animals. \omterm{def:b2:plant-life-cycles:fern}{Fern} strategy wins
in wet, shaded, stable habitats, and in colonising distant or newly
opened ground, where cheap, light, numerous propagules matter more
than provisioning.
\textbf{25.} \omterm{def:b2:plant-life-cycles:fern}{Fern}: $5.1\times 10^{6}$ spores per season for 1.9 adults,
$2.7\times 10^{6}$ spores per adult offspring, \qty{430}{J} each;
pine: $9.6\times 10^{5}$ seeds for 480 adults, \num{2000} seeds per
adult offspring, \qty{100}{kJ} each.
\end{solution}
